\section{Discussion}

The question of whether machines can conduct scientific inquiry has accompanied AI since DENDRAL built explicit rules for generating and eliminating hypotheses from spectral evidence in 1965 \autocite{lindsay1993dendral}. Our data provide a specific answer for current LLM-based agents. 
In workflow-construction domains, agents approach ceiling performance. In hypothesis-driven domains, they do not adapt their reasoning to the demands of the problem: evidence is ignored in 68\% of traces, refutation-driven belief revision occurs in only 26\%, and convergent multi-test evidence is rare. These patterns persist across scaffold architectures, across domains, and under interventions that provide near-complete successful trajectories. They are properties of the base model.

The epistemic process by which a result is reached determines whether the resulting knowledge is justified \autocite{pittphilsci28744,krenn2022understanding, sep-knowledge-analysis}. 
An agent that reaches a correct answer while ignoring contradictory evidence and leaving hypotheses untested produces a result whose reliability on new problems is unknown. Much current engineering effort focuses on scaffolding (prompting strategies, orchestration logic, tool interfaces), but the reasoning patterns we document persist across all scaffold conditions we tested, including extreme interventions. Addressing them will likely require changes to how base models are trained: a slower process than scaffold improvement, and one currently not guided by epistemic criteria. The framework we introduce can supply such guidance: each environment provides a reproducible task, tools, and a scoring function over agent trajectories, sufficient to define training signals on the reasoning process itself.

Human scientists operate within institutional structures that enforce epistemic norms: peer review, replication, reputational consequences. LLM-based agents are subject to no such constraints. Evidence that reliance on AI assistance can erode independent problem-solving \autocite{liu2026persistence} compounds this gap. Characterizing the epistemic behavior of these systems requires explicit measurement. The evaluation framework we introduce does this across tasks and configurations. New environments, models, scaffolds, and metrics can be contributed, allowing the community to track how the epistemic behavior of these systems changes as models evolve. The methods and findings extend beyond scientific agents to any setting where the reasoning process determines the value of the result.

The tools scientists use shape the science they produce \autocite{culkin1967schoolman,hao2026artificial}. LLM-based agents are now among those tools. Our analysis shows that their epistemic behavior does not exhibit the patterns that characterize scientific reasoning. As long as they are evaluated only by the answers they produce, this difference will remain invisible, and it will shape the knowledge they help produce.